\setcounter{section}{0}

\Needspace{5\baselineskip}
\hypertarget{ux4e16ux754cux6a21ux578bux6982ux8ff0}{%
\section{世界模型概述}\label{ux4e16ux754cux6a21ux578bux6982ux8ff0}}

\Needspace{5\baselineskip}
\hypertarget{ux4e00ux4e16ux754cux6a21ux578bux7684ux5b9aux4e49}{%
\subsection{一、世界模型的定义}\label{ux4e00ux4e16ux754cux6a21ux578bux7684ux5b9aux4e49}}

\Needspace{5\baselineskip}
LLM在数字世界已经具备了强大的能力，但对物理世界的理解却并不充分。世界模型是可能的一种解决方案。如果说LLM的核心是``Next token prediction''，那么世界模型的核心就是``Next state prediction''。世界模型定义为给定当前状态和动作，预测下一个（实际上往往可能是未来若干个，这里简化定义）状态的系统：

P(st+1\textbar st,at)

世界模型作为对物理世界的建模，应当能处理多模态（不仅限于视觉，也包括传感器数据、光谱数据等各种多模态数据）输入，生成显式或隐式的世界状态表征，并有效预测在给定动作下的下一状态。

\Needspace{5\baselineskip}
\hypertarget{ux4e8cux4e16ux754cux6a21ux578bux7684ux8d77ux6e90rlux4e2dux7684ux4e16ux754cux6a21ux578b}{%
\subsection{二、世界模型的起源：RL中的世界模型}\label{ux4e8cux4e16ux754cux6a21ux578bux7684ux8d77ux6e90rlux4e2dux7684ux4e16ux754cux6a21ux578b}}

\Needspace{5\baselineskip}
\hypertarget{ux4e00ux6838ux5fc3ux601dux60f3-17}{%
\subsubsection{（一）核心思想}\label{ux4e00ux6838ux5fc3ux601dux60f3-17}}

人类在感知世界时，会根据自身有限的感官信息在脑海中建立对世界的心理模型。我们做出的决定和行动正是基于这个内部模型。受此认知科学理念的启发，本文探讨了为强化学习（RL）环境构建生成式神经网络模型。

研究者将智能体划分为世界模型和控制模型。世界模型学习环境在空间和时间上的压缩表示；控制模型利用从世界模型中提取的特征作为输入，训练一个策略来完成任务。

这种方法的优势在于，它不仅能高效处理信度分配（Credit Assignment）问题，还能允许智能体在其世界模型生成的虚拟``梦境''中进行自我训练，然后再将策略迁移到真实环境中，这也就是基于模型的强化学习。

\Needspace{5\baselineskip}
\hypertarget{ux4e8cux4e16ux754cux6a21ux578bux7684ux67b6ux6784ux548cux8badux7ec3}{%
\subsubsection{（二）世界模型的架构和训练}\label{ux4e8cux4e16ux754cux6a21ux578bux7684ux67b6ux6784ux548cux8badux7ec3}}

\begin{enumerate}
\def\labelenumi{\arabic{enumi}.}
\item
  感知模型：输入t时刻的视觉像素等多模态信息，用变分自编码器（VAE）等将其压缩成一个小型的向量，即隐变量z\_t。
\item
  预测模型：维护隐状态h\_t，运用类似RNN的结构记录历史记忆，以保证长时间内符合物理定律。输入隐状态h\_t、当前状态信息z\_t和动作输出a\_t对未来状态的预测z*\_t+1。可以直接输出采样，或输出未来状态的概率密度函数。根据预测的z*\_t+1和真实下一状态对应的隐变量z\_t+1的差距构造损失函数学习。
\end{enumerate}

\Needspace{5\baselineskip}
\hypertarget{ux4e09ux63a7ux5236ux6a21ux578bux548cux4e16ux754cux6a21ux578bux7684ux534fux540c}{%
\subsubsection{（三）控制模型和世界模型的协同}\label{ux4e09ux63a7ux5236ux6a21ux578bux548cux4e16ux754cux6a21ux578bux7684ux534fux540c}}

\Needspace{5\baselineskip}
控制模型即被通过强化学习训练的策略网络、价值网络等。一般和世界模型有三种协同方法：

\begin{enumerate}
\def\labelenumi{\arabic{enumi}.}
\item
  控制模型的输入不包括世界模型的输出，但在对控制模型输出进行组采样等的时候，结合世界模型的输出。
\item
  将世界模型的输出一并输入控制模型。
\item
  将二者直接合一，保留多头输出。
\end{enumerate}

\Needspace{5\baselineskip}
\hypertarget{ux56dbux8badux7ec3ux5de5ux4f5cux6d41}{%
\subsubsection{（四）训练工作流}\label{ux56dbux8badux7ec3ux5de5ux4f5cux6d41}}

\Needspace{5\baselineskip}
训练的过程是以下步骤的循环：

\begin{enumerate}
\def\labelenumi{\arabic{enumi}.}
\item
  冻结模型权重，智能体与真实环境互动，收集状态-动作对数据。
\item
  运用数据训练感知模型。
\item
  再冻结感知模型，训练预测模型。
\item
  最后冻结感知模型和预测模型，训练控制模型。
\end{enumerate}

\Needspace{5\baselineskip}
\hypertarget{ux4e94ux5bf9ux6297ux6027ux7b56ux7565ux53caux5176ux907fux514d}{%
\subsubsection{（五）对抗性策略及其避免}\label{ux4e94ux5bf9ux6297ux6027ux7b56ux7565ux53caux5176ux907fux514d}}

由于世界模型只是对真实环境的近似概率模型，控制器在优化过程中可能会发现并利用这些不完美的漏洞，产生``对抗性策略''（例如以一种特定的移动方式让生成的火球凭空熄灭）。适当提高预测模型的温度参数来增加虚拟环境的不确定性和随机性，可以有效防止控制器利用这些模型缺陷。这使得智能体学会了更稳健的策略，从而在真实环境中的表现更好。

\Needspace{5\baselineskip}
\hypertarget{ux4e09ux4e16ux754cux6a21ux578bux4e0eux5177ux8eabux667aux80fdux7684ux5173ux7cfb}{%
\subsection{三、世界模型与具身智能的关系}\label{ux4e09ux4e16ux754cux6a21ux578bux4e0eux5177ux8eabux667aux80fdux7684ux5173ux7cfb}}

\Needspace{5\baselineskip}
\hypertarget{ux4e00ux4e16ux754cux6a21ux578bux9700ux8981ux5177ux8eabux4ea4ux4e92}{%
\subsubsection{（一）世界模型需要具身交互}\label{ux4e00ux4e16ux754cux6a21ux578bux9700ux8981ux5177ux8eabux4ea4ux4e92}}

想要让模型真正理解物理世界的规律，只看视频是不够的。视频数据存在一些显著的问题。

\begin{enumerate}
\def\labelenumi{\arabic{enumi}.}
\item
  缺乏动作条件：视频数据只有状态的变化，却没有动作条件。真正的物理理解常来自干预：``如果我推一下，会不会倒？''``如果我加热，会不会变色？''这些往往依赖于动作条件数据。如果没有动作条件，那么就只是在训练视频生成模型。
\item
  缺乏触觉和力学数据：目前的世界模型训练仍以视觉数据为主。但如果希望世界模型真正理解物理规律，那么模型不仅要``看懂''世界，还能懂得摩擦、硬度等物理性质，这需要触觉和力学数据，以免让模型生成``视觉上合理、物理上错误''的预测，如物体没有支撑力、碰撞没有动量传递等。
\end{enumerate}

\Needspace{5\baselineskip}
\hypertarget{ux4e8cux4e16ux754cux6a21ux578bux5bf9ux5177ux8eabux667aux80fdux53d1ux5c55ux81f3ux5173ux91cdux8981}{%
\subsubsection{（二）世界模型对具身智能发展至关重要}\label{ux4e8cux4e16ux754cux6a21ux578bux5bf9ux5177ux8eabux667aux80fdux53d1ux5c55ux81f3ux5173ux91cdux8981}}

\begin{enumerate}
\def\labelenumi{\arabic{enumi}.}
\item
  世界模型是World-Action Model的重要组成部分，与策略模型结合构成具身基础模型。
\item
\Needspace{5\baselineskip}
  世界模型可以用作训练和评测环境。真实世界中的试错不仅效率极低，而且容易造成昂贵的硬件损耗。具备世界模型的具身智能体可以在内部生成的空间中运行，从而将世界模型作为训练和评测环境。实操时，需要注意几个点：
\end{enumerate}

（1）由于世界状态预测（视频生成）需要耗费的算力显著大于动作预测需要的算力，因此我们不必逐帧预测，而是采用块级预测。

（2）机器人往往有多个摄像头，这些摄像头所获得的图像信息必须满足三维一致性，即共同构成一个一致的三维世界。一种方法是可以每次先生成其中一个视角，再把该视角的图像放入上下文，以此作为条件生成下一个视角，以此类推。

\begin{enumerate}
\def\labelenumi{\arabic{enumi}.}
\setcounter{enumi}{2}
\tightlist
\item
  世界模型可以作为合成数据引擎。世界模型可以生成大量带有物理反馈交互逻辑的合成轨迹，以极低的边际成本反哺动作策略的训练。
\end{enumerate}

\FloatBarrier
\Needspace{5\baselineskip}
\hypertarget{ux53c2ux8003ux6587ux732e-154}{%
\subsection{参考文献}\label{ux53c2ux8003ux6587ux732e-154}}

\vspace{0.25em}
\begin{itemize}
\tightlist
\item
  Ha, D., \& Schmidhuber, J. (2018). \href{https://arxiv.org/abs/1803.10122}{World Models}. arXiv:1803.10122.
\item
  Sutton, R. S. (1991). \href{https://dl.acm.org/doi/10.1145/122344.122377}{Dyna, an Integrated Architecture for Learning, Planning, and Reacting}. ACM SIGART Bulletin.
\end{itemize}

\clearpage
